\section{Prime--squarefree logarithmic sampling}
\label{sec:tpc51-arithmetic-sampling}

We now verify the finite-Gram hypothesis on the row universe actually
used before terminal prime summation.  Only standard fixed-modulus
prime number theory and M\"obius inversion are needed.

Put \(T=\log 2\).  Fix positive integers \(q,H\), a nonempty set
\(\cA\subset(\Z/q\Z)^2\), and \(e\in\{0,1\}\).  We assume that every
prime residue occurring in \(\cA\) is reduced modulo \(q\), and retain
only those squarefree residue classes having positive local density
subject to \((d,H)=1\).  At least one retained pair is required.  Define
\begin{align}
 \cR_X(\cA)
 &:=\bigl\{(\ell,d):L<\ell\leq2L,\ \ell\ {\rm prime},
       \ D<d\leq2D,\ \mu^2(d)=1,\ (d,H)=1,\notag\\[-2pt]
 &\hspace{47mm}(\ell,d)\bmod q\in\cA\bigr\},
 \label{eq:tpc51-row-universe}\\
 w_{\ell,d}&:=(\log\ell)^{2e},
 \qquad
 E_X:=\sum_{(\ell,d)\in\cR_X(\cA)}w_{\ell,d}.
 \label{eq:tpc51-row-base-energy}
\end{align}
We suppose \(X^c\leq L,D\leq X^C\) for fixed \(c,C>0\).  The normalized
logarithmic empirical measure is
\begin{equation}
 \nu_X:=\frac1{E_X}
 \sum_{(\ell,d)\in\cR_X(\cA)}
 w_{\ell,d}\,
 \delta_{(\log(\ell/L),\log(d/D))}.
 \label{eq:tpc51-empirical-measure}
\end{equation}

For the quantitative comparison it is convenient not to discard the
harmless finite-\(L\) logarithmic density.  For all sufficiently large
\(L\), let
\begin{equation}
 Z_{L,e}:=\int_0^T e^u(\log L+u)^{2e-1}\,du
 \label{eq:tpc51-prime-reference-normalizer}
\end{equation}
and define the probability measure
\begin{equation}
 d\nu_{L,e}^{\rm ref}(x,y)
 :=\frac{e^x(\log L+x)^{2e-1}}{Z_{L,e}}\,dx\,
   e^y\,dy,
 \qquad (x,y)\in[0,T]^2.
 \label{eq:tpc51-finite-reference-measure}
\end{equation}
The second factor already has mass one because
\(\int_0^T e^y\,dy=1\).  Notice that the exponent \(2e-1\) equals
\(-1\) when the prime row is unweighted; this is the usual density
\(dt/\log t\) of primes and is the reason for retaining the finite-\(L\)
reference measure.

\begin{theorem}[Arithmetic logarithmic sampling]
\label{thm:tpc51-row-equidistribution}
For every continuous \(F:[0,T]^2\to\C\),
\begin{equation}
 \int F\,d\nu_X
 \longrightarrow
 \int_{[0,T]^2}F(x,y)e^{x+y}\,dx\,dy.
 \label{eq:tpc51-weak-row-limit}
\end{equation}
More quantitatively, for every fixed \(A>0\) and every
\(0<\delta<1/2\), uniformly when
\(|\xi|+|\eta|\leq(\log X)^A\),
\begin{align}
 &\left|
 \int e^{i(\xi x+\eta y)}\,d\nu_X
 -\int e^{i(\xi x+\eta y)}\,d\nu_{L,e}^{\rm ref}
 \right|
 \notag\\
 &\hspace{18mm}\leq
 C_{A,\delta,q,H,\cA}
 (1+|\xi|+|\eta|)^{C_1}
 \left(e^{-c_1\sqrt{\log L}}+D^{-1/2+\delta}\right),
 \label{eq:tpc51-fourier-moment-error}
\end{align}
where \(c_1>0\) and \(C_1<\infty\) are independent of
\(X,\xi,\eta\).
\end{theorem}

\begin{proof}
Split \(\cA\) into its finitely many residue pairs.  For a reduced
prime class \(a\bmod q\), the fixed-modulus prime number theorem and
partial summation give, uniformly for
\(|\xi|\leq(\log X)^A\),
\begin{align}
 &\sum_{\substack{L<\ell\leq2L\\
                   \ell\ {\rm prime}\\
                   \ell\equiv a\pmod q}}
 (\log\ell)^{2e}e^{i\xi\log(\ell/L)}
 \notag\\
 &\quad=
 \frac{L}{\varphi(q)}
 \int_0^T e^x(\log L+x)^{2e-1}e^{i\xi x}\,dx
 +O_{A,q}\!\left(
 L(\log L)^{2e-1}(1+|\xi|)^{C_1}
 e^{-c_1\sqrt{\log L}}\right).
 \label{eq:tpc51-prime-fourier-asymptotic}
\end{align}
The same formula at \(\xi=0\) supplies its normalizing mass.  The use
of the finite-\(L\) reference density in
\eqref{eq:tpc51-finite-reference-measure} removes any artificial
\(1/\log L\) error from this comparison.  Any fixed extra power of
\(\log L\) arising in partial summation is absorbed by decreasing
\(c_1\), without changing the stated form of the error.

For a locally admissible squarefree class \(b\bmod q\), use
\(\mu^2(d)=\sum_{r^2\mid d}\mu(r)\) and inclusion--exclusion for
\((d,H)=1\).  For each \(r\leq\sqrt{2D}\) and each fixed divisor
\(s\mid H\), the remaining integer is counted in either no class or one
class modulo \(\operatorname{lcm}(q,s,r^2)\).  Summing the resulting
\(O(1)\) endpoint errors over \(r\), extending the absolutely
convergent main-term sum to infinity, and applying partial summation to
\(e^{i\eta\log(d/D)}\) give
\begin{align}
 &\sum_{\substack{D<d\leq2D\\d\equiv b\pmod q\\
       \mu^2(d)=1,\ (d,H)=1}}
 e^{i\eta\log(d/D)}
 \notag\\
 &\quad=
 \mathfrak c_{b,q,H}D
 \int_0^T e^y e^{i\eta y}\,dy
 +O_{A,\delta,q,H}\!\left(
 D^{1/2+\delta}(1+|\eta|)^{C_1}\right),
 \label{eq:tpc51-squarefree-fourier-asymptotic}
\end{align}
uniformly for \(|\eta|\leq(\log X)^A\), with
\(\mathfrak c_{b,q,H}>0\) precisely for the retained classes.  This is
the usual squarefree count in a fixed progression with a smooth
unit-modulus test factor; no distribution theorem for growing moduli is
being used.

Multiply \eqref{eq:tpc51-prime-fourier-asymptotic} and
\eqref{eq:tpc51-squarefree-fourier-asymptotic}, sum over the fixed set
\(\cA\), and divide by the same expression at
\((\xi,\eta)=(0,0)\).  The residue-dependent main-term constants form
one positive finite sum, while the \(x\)- and \(y\)-profiles are common
to every retained residue pair.  The normalized main term is therefore
exactly \(\nu_{L,e}^{\rm ref}\), and the denominator is bounded away
from zero after its natural normalization.  This proves
\eqref{eq:tpc51-fourier-moment-error}.

For fixed frequencies the normalized prime reference density converges
uniformly to \(e^x\).  Hence
\(\nu_{L,e}^{\rm ref}\) converges weakly to the probability measure
\(e^{x+y}\,dx\,dy\).  The quantitative estimate first proves
\eqref{eq:tpc51-weak-row-limit} for finite exponential sums with
arbitrary real frequencies.  This self-adjoint algebra contains the
constants and separates points of the closed square, so the
Stone--Weierstrass theorem gives the result for every continuous \(F\).
\end{proof}

\subsection{A polylogarithmic log-Fourier bank}

For an integer \(K\geq0\), let
\(\Lambda_K=\{-K,\ldots,K\}^2\) and
\(R_K=|\Lambda_K|=(2K+1)^2\).  For coefficients
\(a=(a_k)_{k\in\Lambda_K}\), put
\begin{equation}
 P_a(x,y):=\sum_{k=(k_1,k_2)\in\Lambda_K}
 a_k\exp\!\left(\frac{2\pi i}{T}(k_1x+k_2y)\right).
 \label{eq:tpc51-log-fourier-polynomial}
\end{equation}

\begin{theorem}[Polylogarithmic sampled Riesz bounds]
\label{thm:tpc51-polylog-fourier-frame}
Fix \(A>0\).  If \(K\leq(\log X)^A\), then for all sufficiently large
\(X\), uniformly in \(a\),
\begin{equation}
 c_A\sum_{k\in\Lambda_K}|a_k|^2
 \leq \int|P_a|^2\,d\nu_X
 \leq C_A\sum_{k\in\Lambda_K}|a_k|^2,
 \label{eq:tpc51-unmasked-fourier-frame}
\end{equation}
where \(c_A>0\) and \(C_A<\infty\) are independent of \(K,X\).
The same statement holds for coefficients in an arbitrary Hilbert
space, with \(|P_a|^2\) replaced by its squared Hilbert norm.
\end{theorem}

\begin{proof}
For all sufficiently large \(L\), the density in
\eqref{eq:tpc51-finite-reference-measure} is bounded above and below by
positive absolute constants on \([0,T]^2\).  Parseval on the square
therefore gives
\begin{equation}
 c_0\sum_k|a_k|^2
 \leq\int|P_a|^2\,d\nu_{L,e}^{\rm ref}
 \leq C_0\sum_k|a_k|^2.
 \label{eq:tpc51-reference-frame}
\end{equation}
Expanding the quadratic form produces frequencies whose coordinates
have absolute value at most \(4\pi K/T\).  By
\eqref{eq:tpc51-fourier-moment-error}, applied with any fixed exponent
larger than \(A\) so that the harmless constant \(8\pi/T\) is absorbed
for large \(X\), every entry of the empirical
Gram differs from the reference Gram by at most
\begin{equation}
 \rho_X(K)\ll_{A,\delta}
 (1+K)^{C_1}
 \left(e^{-c_1\sqrt{\log L}}+D^{-1/2+\delta}\right).
 \label{eq:tpc51-bank-entry-error}
\end{equation}
The Schur bound, or simply
\((\sum|a_k|)^2\leq R_K\sum|a_k|^2\), shows that the quadratic-form
error is at most
\begin{equation}
 R_K\rho_X(K)\sum_k|a_k|^2.
 \label{eq:tpc51-bank-operator-error}
\end{equation}
Because \(R_K\ll(1+K)^2\), \(K\leq(\log X)^A\), and
\(D\geq X^c\), this is \(o(1)\sum_k|a_k|^2\).  Absorb the error into
\eqref{eq:tpc51-reference-frame}.  For Hilbert-valued coefficients the
scalar Gram is tensored with the identity, so the same spectral bounds
apply.
\end{proof}

\begin{corollary}[Fixed physical families]
\label{cor:tpc51-fixed-physical-family}
Let \(W_1,\ldots,W_R\in C([0,T]^2)\) be fixed.  Quotient their span by
the exact relations in \(L^2(e^{x+y}\,dx\,dy)\), and choose fixed
representatives whose nonzero quotient classes are linearly independent.
Then the normalized sampled Gram of this reduced family on
\(\cR_X(\cA)\) has its least eigenvalue bounded away from zero.  In
particular, its quotient-frame condition number is \(O(1)\).
\end{corollary}

\begin{proof}
Apply \cref{thm:tpc51-row-equidistribution} to each of the finitely many
continuous products \(W_r\overline{W_s}\).  The resulting sampled Gram
converges in operator norm to the positive-definite limiting Gram, so
\cref{thm:tpc51-limit-gram} applies.
\end{proof}

The corollary is the direct tool for the finite collection of fixed
physical cutoff shapes.  The growing Fourier bank is a stronger stress
test and a convenient finite recellization device; it is not being
identified with an arbitrary continuous Mellin coefficient space.
